% Recovered verbatim from chapters_11_15_with_part/chapters_11_15_with_part.tex
% Canonical chapter 11
\chapter{Phase Cosmology and Background Solutions}

\section{Homogeneous and Isotropic Ansatz}

We consider a spatially homogeneous and isotropic cosmology with metric
\[
  ds^2 = -dt^2 + a(t)^2 \left( \frac{dr^2}{1-kr^2} + r^2 d\Omega^2 \right),
\]
where $a(t)$ is the scale factor and $k=0,\pm 1$ the spatial curvature index.

The phase field and chromatic fields are taken to be spatially homogeneous:
\[
  \tau = \tau(t), \qquad \chi_a = \chi_a(t),
\]
while the ER network contributes an effective homogeneous energy component $\rho_{ER}(t)$.

\section{Effective Friedmann Equations}

From the MK--43 Ultra action we obtain modified Friedmann equations of the form
\begin{align}
  H^2 &=
  \frac{8\pi G}{3}\left(\rho_\tau + \rho_\chi + \rho_{ER} + \rho_{\mathrm{m}}\right)
  - \frac{k}{a^2},
  \\
  \frac{\ddot{a}}{a} &=
  -\frac{4\pi G}{3}\left(
    \rho_\tau + \rho_\chi + \rho_{ER} + \rho_{\mathrm{m}}
    + 3p_\tau + 3p_\chi + 3p_{ER} + 3p_{\mathrm{m}}
  \right),
\end{align}
where $\rho_\tau$, $\rho_\chi$, and $\rho_{ER}$ denote energy densities in the phase, chromatic, and ER sectors, respectively, and $p_\tau$, $p_\chi$, $p_{ER}$ are the associated pressures.

\section{Role of the Phase Field in Expansion}

The phase field contributes an effective equation of state $w_\tau = p_\tau/\rho_\tau$ which can interpolate between stiff, radiation-like, or dark-energy-like behavior depending on the potential $V(\tau)$ and the magnitude of $\dot{\tau}$.
In particular, quasi-de Sitter phases can arise from regions where $V(\tau)$ dominates and $\dot{\tau}$ is small.

